\chapter{The Build Stage}

The procedure for creating the source distribution.

\section{Final checks}

\begin{itemize}

\item Check plug-in dependencies in all \texttt{configure.*}, in case they have
  changed: if you don't know them, ask to plug-in developers to verify them.

\todo{should be done when the plugin is modified}

\item Check the headers and corresponding licenses files for new files
  \todo{should be done when the plugin is modified}

%  (\texttt{make check-headers} is your friend, but it is done by CI anyway).


\item Check the contents of \texttt{INSTALL.md} \todo{Should always be up to date}
\item Check the contents of \texttt{README.md} \todo{Should always be up to date}
\item Check the contents of \texttt{Changelog}
  \todo{Should always be up to date}
\item Carefully read the output of the configure to be sure that there are no
  buggy messages. Make sure you have required OPAM packages installed in your machine,
  e.g. alt-ergo.
\end{itemize}

\section{Update the Sources}\label{sec:update-sources}

There are many administrative steps, coordinated by the release manager.
\begin{enumerate}
\item Ensure that you have a clone of \texttt{e-acsl}
  (\texttt{git@git.frama-c.com:frama-c/e-acsl}) in the \texttt{src/plugins}
  directory
\item Set \texttt{VERSION} file (for beta releases, add suffix \texttt{-beta}).
\item Update file \texttt{ALL\_VERSIONS} (nothing to do for beta releases). \todo{This section must be put somewhere else in the release manual}
\item Update files
  \texttt{src/plugins/wp/Changelog} and
  \texttt{src/plugins/e-acsl/doc/Changelog},
  to add the header corresponding to the
  new version. % For the final release, use the script
  % \texttt{doc/changelog/generate} to check that the HTML page can be built,
  % and check its contents.
\item Update the Frama-C's authors list in files
  \texttt{src/plugins/gui/help\_manager.ml} and \texttt{opam/opam}.
\item Recompile and test to check that all is fine.
\item Check the documentation as per Section~\ref{update_doc}
\item update the \texttt{opam} directory: update the version name and
  version number in the fields \texttt{version}, \texttt{doc} and
  \texttt{depends} in \texttt{opam/opam}
\item Create the last commit
\item Create the tag using \texttt{git tag \$(cat VERSION)} and push it (e.g. via \texttt{git push origin \$(cat VERSION)})
\end{enumerate}

\section{Update the Documentation}\label{update_doc}

\subsection{Manuals}

Most manuals depend on the \texttt{VERSION} file, so make sure to regenerate
them when this file changes.
Also, most of the manuals include an appendix with the list of
changes that have happened, divided by version of Frama-C. Make sure that the
first section has the correct release name.

\todo{Mention the other things that change from one release to the
  other, such as the ACSL version number}

\todo{The release manager should create fresh Changes sections in \texttt{master}. But When?: at fork time, or when the release is ready?}


The official \FramaC documents relative to the kernel and ACSL to be released
must be recompiled in order to take into account the \texttt{VERSION}
file:

\begin{center}
  \begin{tabular}{lll}
    \hline
    \textsf{User Manual} &
    \texttt{doc/userman} &
    \expertise{Julien} \\
    \textsf{Developer Manual} &
    \texttt{doc/developer} &
    \expertise{Julien} \\
    \textsf{ACSL Implementation}  &
    \texttt{doc/acsl} (from Github) &
    \expertise{Virgile} \\
    % \textsf{man page} &
    % \texttt{man/frama-c.1} &
    % \expertise{Virgile} \\
    \hline
  \end{tabular}
\end{center}

Documents relative to the plug-ins can be found at the same place. This list
isn't exhaustive:

\begin{center}
  \begin{tabular}{lll}
    \hline
    \textsf{RTE}  &
    \texttt{doc/rte} &
    \expertise{Julien} \\
    \textsf{Aoraï}  &
    \texttt{doc/aorai} &
    \expertise{Virgile} \\
    \textsf{Metrics}  &
    \texttt{doc/metrics} &
    \expertise{Richard} \\
    \textsf{Value Analysis}  &
    \texttt{doc/value} &
    \expertise{Boris} \\
    \textsf{WP}  &
    \texttt{src/plugins/wp/doc/manual/} &
    \expertise{Loïc} \\
  \end{tabular}
\end{center}

Clone the \texttt{acsl} manual
(\url{git@github.com:acsl-language/acsl.git}) in
\texttt{doc/acsl} in order to generate the ACSL reference and implementation
manuals.

% manuals (which uses git-annex) are now deprecated in favor of using the
% Github wiki
%Clone the \texttt{manuals} repository
%(\url{git@git.frama-c.com:frama-c/manuals.git}) in
%\texttt{doc/manuals} in order to update the manuals (required by the
%\texttt{install} target below).

In each directory:
\begin{shell}
make
make install
\end{shell}

The result is installed in doc/manuals/.

Finally, E-ACSL manuals must also be generated and installed: in both
directories \texttt{src/plugins/e-acsl/doc/refman}\footnote{If you followed
instructions from section~\ref{sec:update-sources}, you have already cloned
the E-ACSL repository under \texttt{src/plugins/e-acsl}.
Otherwise, it's time to do it}
and
\texttt{src/plugins/e-acsl/doc/userman}, do
\begin{shell}
make
make install
\end{shell}


The script \texttt{doc/build-manuals.sh} performs this for all the current
manuals (including E-ACSL).

\begin{shell}
./doc/build-manuals.sh
\end{shell}

% No longer use git annex (as recommended by VP)
%\paragraph{Getting the manuals on the web (ultimately)}
%
% In \texttt{doc/manuals} after installing all the manuals:
%\begin{shell}
%git annex add .
%git annex copy . --to origin
%git commit -am "manuals for release ...."
%git push
%\end{shell}

\subsection{API documentation}

To update the \textsf{API} documentation, you must update the
\texttt{@since} tags. This involves the following script (\texttt{sed -i} must
work for this purpose):
\begin{shell}
./bin/update_api_doc.sh <NEXT>
\end{shell}
\verb+<NEXT>+ is the name of the release, i.e. the contents of VERSION (without any \verb+rc+ suffix).  

Check that no \verb|+dev| suffix remains inside comments:

\begin{shell}
git grep +dev src
\end{shell}

Test by recompiling Frama-C and rebuilding the API documentation:
\begin{shell}
make -j byte gui-byte
make doc # or just doc-kernel
\end{shell}

\todo{What is this step:
check consistency of API documentation. (the plug-in development guide
  (\texttt{make check-devguide}) doesn't work anymore)}

\section{Build the Source Distribution}
\label{sec:build-source-distr}

This steps creates the tarball of Frama-C, the tarball of the API and
copy them to the website. It also copies the manuals.
You need to have installed \texttt{git-lfs}, otherwise the build
script will not correctly use lfs when putting large files in the
website git.

\expertise{release manager} Use the script
\texttt{build-src-distrib.sh} for this purpose (\texttt{bash version
  4} required) after cloning the following repositories
\begin{itemize}
\item \texttt{Frama-C-snapshot} (\texttt{git@github.com:Frama-C/Frama-C-snapshot}),
\item \texttt{Frama-C-snapshot.wiki} (\texttt{git@github.com:Frama-C/Frama-C-snapshot.wiki}), and
\item \texttt{website} (\texttt{git@git.frama-c.com:frama-c/website})
\end{itemize}
in the root directory of Frama-C. The last one is only needed when creating a
final release. Release candidates are only put on github and require only
\texttt{Frama-C-snapshot} and its associated wiki. Branch on
\texttt{Frama-C-snapshot} should be \texttt{master} for a final release and
\texttt{release-candidates} for a release candidate. For \texttt{website},
a new branch should be created over \texttt{online}.

Finally, ensure that locale \verb+en_US.UTF-8+ is available on your system,
as it is used by the script to ensure a proper encoding of the generated files.

Now, run the script:
\begin{shell}
./bin/build-src-distrib.sh
\end{shell}

\section{Testing the Distribution}

Please check that the distribution (sources and API) is OK:
\begin{enumerate}
\item Inspect the contents of the \texttt{tar.gz}.
  This is done in two steps:
  \begin{enumerate}
    \item run \texttt{doc/release/checktar.sh}. Fix issues and/or update
      the script if needed.
    \item On a pristine clone of the Github Frama-C-snapshot repository,
  compare the untared directory and the Github clone using e.g. \texttt{meld},
  to see if there are suspicious newly-added files or directories (indicating
  e.g. that the .tar.gz was created on a non-clean environment).
  Note that comparing file contents is irrelevant at this stage, since
  there are likely too many differences between this and the last release.
  However, check for post-release commits (namely to \texttt{opam/opam}) and
  see if there are changes that were forgotten to be ported back to Frama-C's git.
\end{enumerate}
\item check for possible path/environment leaks in the \texttt{tar.gz}:
  grep for the path where your Frama-C git directory is, e.g.
  \texttt{/home/user/git/frama-c}. It should appear nowhere in the archive.
\item check that \texttt{./configure \&\& make \&\& sudo make install} works fine;
\item test the installed binaries (especially the GUI). (On Linux, OCI tests
  everything but the GUI);
\item redo the two steps above on Windows \expertise{André}, Mac OS X
  \expertise{Loïc/Thibaud};
\item have a look at the documentation of the API. \todo{What, precisely}
\end{enumerate}


%%%Local Variables:
%%%TeX-master: "release"
%%%End:
